
El presente Manual de Control de Mantenimiento (MCM) de SKYFILMS ha sido elaborado con el propósito de establecer las políticas, procedimientos, responsabilidades y estándares que rigen la gestión del mantenimiento de los Sistemas de Aeronaves No Tripuladas (UAS) operados por la organización.\\

Este manual constituye el documento oficial de referencia técnica y administrativa para asegurar que todas las actividades de mantenimiento preventivo, correctivo, de inspección y control técnico de los equipos se realicen conforme a los Reglamentos Aeronáuticos de Colombia (RAC 100) y a las directrices emitidas por la Unidad Administrativa Especial de Aeronáutica Civil (Aerocivil).\\

El MCM busca garantizar que la flota de UAS de SKYFILMS se mantenga en condiciones de aeronavegabilidad continuada, preservando en todo momento la seguridad operacional, la confiabilidad técnica y la eficiencia de las operaciones aéreas realizadas con fines de producción audiovisual, fotografía aérea y servicios relacionados.\\

Asimismo, este documento establece los mecanismos de control documental, trazabilidad y verificación técnica, aplicables tanto al mantenimiento ejecutado directamente por SKYFILMS como al realizado por proveedores externos aceptados, bajo los lineamientos de calidad y seguridad definidos por la organización.\\

\section{Alcance}

El presente Manual aplica a todas las actividades de mantenimiento, inspección, reparación, reemplazo, control técnico y registro de los UAS pertenecientes a SKYFILMS, incluyendo sus componentes, baterías, equipos de control remoto, sensores y accesorios asociados.\\

Su alcance comprende:

\begin{itemize}
    \item Los procesos internos de mantenimiento realizados en las instalaciones de SKYFILMS.

    \item El mantenimiento tercerizado, efectuado por organizaciones o técnicos externos previamente evaluados y aceptados.

    \item Las acciones de control, documentación y seguimiento a las condiciones técnicas y de aeronavegabilidad de los UAS, durante todas las fases operacionales.

    \item La implementación de programas de mantenimiento preventivo, inspecciones periódicas, control de herramientas, manejo de baterías, almacenamiento y disposición final de componentes.
\end{itemize}

El MCM aplica al personal técnico, operativo y administrativo que participa directa o indirectamente en los procesos de mantenimiento de UAS, y es de cumplimiento obligatorio para todos los colaboradores y contratistas que realicen actividades bajo el certificado de operación de SKYFILMS.

\section{Marco Normativo y Referencias}
Este Manual se fundamenta en la normatividad aeronáutica vigente en la República de Colombia, particularmente en:

\begin{itemize}
    \item \textbf{RAC 100} – Operación de Sistemas de Aeronaves No Tripuladas (UAS).

    \item \textbf{MAUT-5.0-22-011} – Guía para obtener el certificado de explotador UAS.

    \item \textbf{MAUT-5.0-12-097FO} – Lista de verificación para revisar el Manual de Control de Mantenimiento.

    \item Manual del fabricante de cada modelo de UAS operado por SKYFILMS.
\end{itemize}

El cumplimiento de este marco normativo es obligatorio y verificable ante la Aerocivil, y todo el personal técnico de SKYFILMS deberá conocer, aplicar y mantener actualizada la información contenida en el presente documento.

